%%%%%%%%%%%%%%%%%%%%%%%%%%%%%%%%%%%%%%%%%%%%%%%%%%%%%%%%%%%%%
%                                                           %
% Cette version repart de 03_Draft\rolefe_current-short.tex %
%                                                           %
%%%%%%%%%%%%%%%%%%%%%%%%%%%%%%%%%%%%%%%%%%%%%%%%%%%%%%%%%%%%%

% NB: Je supprime toutes les (grosses) parties qui étaient en commentaire dans la version précédente pour qu'on y voit plus clair.


\subsection{Empirical specification}\label{sec:empirical}

We analyze the cost of job displacement by estimating the following equation:
\begin{equation}\label{eq:main_estimating_equation}
Y_{it} = \sum\limits_{d=-2}^{6} \beta_{d} \times {\mathbbm{1}}\{t^D_i+d=t\} \times \mathbbm{1}\{D_i=1\} + \sum\limits_{d=-3}^{6} \delta_{d} \times {\mathbbm{1}}\{t^D_i+d=t\} + \alpha_i + \delta_t + \epsilon_i,
\end{equation}

where $Y_{it}$ is an outcome of interest measured in year $t$ for individual $i$ employed in firm $J(it)$, displaced in year $t^D_i$. $\beta_{d}$ measures the change in $Y$ for displaced workers (indicated by $D=1$) compared to the control group, $d$ periods after displacement. The model includes worker ($\alpha_i$) and year ($\delta_t$) fixed effects. In these estimations, we take $t^D-2$ as the reference year, as outcomes may already be affected in the year prior to displacement ($t^D-1$). Standard errors are clustered at the individual level.\footnote{We use a matching procedure with replacement: non-displaced individuals may enter the control group multiple times if they are the best match for several displaced workers. In this case they will appear under different ``worker'' identities. We cluster standard errors at the true ``individual'' level however to account for the dependence between observations that this may introduce. }

\subsection{Total earnings and employment}\label{sec:earnings_employment} %and hours worked
We first establish the impact of displacement on earning and employment, including adjustment at the extensive margin. Table \ref{tab:evstudytable_body} reports the coefficients $\beta_{d}$ from different estimations of equation (\ref{eq:main_estimating_equation}).\footnote{Throughout this section, the headline coefficients on hours and earnings are reported in log points; these should be read as approximate level percentages only when small. The largest coefficients we report ($-0.939$ on log earnings and $-0.961$ on log hours at $t^D$ -> \textcolor{red}{-0.932 / -0.953} \textcolor{orange}{[OK: Table 2, cols. 3--4 at $t^D$]}) correspond to level drops of $1-\exp(-0.939) = 60.9\%$ and $1-\exp(-0.961) = 61.7\%$ \textcolor{orange}{[$1-\exp(-0.932)=60.6\%$ and $1-\exp(-0.953)=61.4\%$]} respectively, not to the near-elimination of earnings that the log values might suggest if read linearly. The level specification in column~(2) of Table~\ref{tab:evstudytable_body} provides a direct check: the earnings effect at $t^D$ is $-14{,}062$\EUR{} -> \textcolor{red}{13,731} \textcolor{orange}{[OK: Table 2, col. 2 at $t^D$]} on a control-group mean of $25{,}425$\EUR{} -> \textcolor{red}{25,710} \textcolor{orange}{[control mean not in any table; 13,731/25,710 = 53.4\% is arithmetically right]} at the same horizon, a $55.3\%$ -> \textcolor{red}{53.4\%} drop, in line with the log estimate.} In columns (1) and (2), the outcome variable is employment status and overall earnings, respectively. Because outcomes are measured at the annual level and displacement occurs within the year, estimates at $t^D$ reflect a mix of pre- and post-displacement outcomes and therefore understate the full effect, which is more accurately captured at $t^D+1$. The results show the well-known losses following displacement: employment probability falls by about 32 percentage points at $t^D+1$ (relative to the control group), then progressively attenuates over time to about $-14$ percentage points after five years. This fall drives a decline in annual earnings that culminates (\textcolor{red}{false: it doesn't culminate, so either we go to 0 or change the phrasing}) \textcolor{orange}{[Table 2, col. 2: $-13{,}731$ at $t^D$, $-13{,}687$ at $t^D+1$, $-8{,}690$ at $t^D+5$; proposed: ``This fall drives a decline in annual earnings of about 13,700 euros in the year following displacement (55\% of the baseline annual earnings of Table 1), which attenuates to about 8,700 euros (35\%) at $t^D+5$.'']} one year following displacement at about \euro 13,900 -> \textcolor{red}{13,687} \textcolor{orange}{[OK]} (39\% -> \textcolor{red}{55\%} \textcolor{orange}{[OK: 13,687/24,848]} of the baseline annual earnings of Table \ref{tab:desc:AllLE}) and attenuates to about \euro 8,700 -> \textcolor{red}{8,670} \textcolor{orange}{[8,690: Table 2, col. 2 at $t^D+5$]} (25\% -> \textcolor{red}{35\%} \textcolor{orange}{[OK: 8,690/24,848]} of baseline earnings) at $t^D+5$. Appendix Figure \ref{fig:aveffect_emp_earn} display those results and associated pre-trends. Figures \ref{fig:raw_mean_earnings} and \ref{fig:raw_mean_hours} present the raw means for each group. 

\subsection{Decomposition of earnings losses}\label{sec:decomposition_losses}
We now decompose earnings losses into hours-worked losses and hourly-wage losses, and further decompose hourly-wage losses, within the AKM framework, into a firm wage-premia component and a residual component. As is standard in the literature, we focus on re-employed workers. This conditioning matters for interpretation: displaced workers are 22 percentage points less likely to be employed than their controls at $t^D+2$, and still 14 points less likely at $t^D+5$ (Table~\ref{tab:evstudytable_body}, column 1), so the estimates below are computed on a selected sample, whose composition changes over time as workers who return to work later progressively enter it. The hourly-wage loss is nevertheless stable between $t^D+2$ ($-0.106$) and $t^D+5$ ($-0.101$), while the loss in firm premia grows from $-0.029$ to $-0.037$, which suggests that later entrants earn about as much per hour as earlier ones, but in firms paying lower premia. Figure \ref{fig:fe_current_event} displays the various elements of this decomposition: losses in total earnings, hours worked, hourly-wage and firm premia. Earnings and hours (in log) drop abruptly and slowly recover. The hours margin, even conditional on being re-employed, remain very strong. This highlights the importance of working with data on hours to assess properly the relative contribution of labor supply and wage rate to the overall earning response.  Interestingly, the effect on hourly wage stabilizes in the medium run at around -0.10, so its contribution to overall loss in earnings grows over time.\footnote{We also observe a counter-intuitive yet familiar bump immediately after the layoff, driven by severance payment, delayed pay and potential measurement error in hours worked.} This finding is in line with recent research on displacement in other countries or over different periods, such as \cite{Lachowska2020} or \cite{Gulyas2019}. While \cite{Schmieder2020} show evidence of a recovery in \textit{daily} wage rate, they also find that the contribution of wage rate decline to overall earnings loss grows over time. (\textcolor{red}{main decomposition isn't same sample, we exported it (it's called sample\_nonmiss in main results) but we never used it beside I think in ref reports?})

\paragraph{Contribution of employer premia to losses in wage rates.} We assess the importance of the loss in firm-specific wage premia using the firm fixed effect as a dependent variable in our main model (equation \ref{eq:main_estimating_equation}). Coefficients are reported in Figure \ref{fig:fe_current_event}. In line with most of the literature, we find that displaced workers re-employ in (relatively) worse-paying firms. By $t^D+5$ the average loss in premium reaches {-0.072} -> \textcolor{red}{-0.037 (note: 7.2 is from the first paper version, no?)} \textcolor{orange}{[$-0.037$ OK: Table 2, col. 6 at $t^D+5$; the $-0.072$ was already stale in R1, whose table had $-0.042$]}. 


To assess the contribution of firm wage premia to losses in hourly wages, \cite{Lachowska2020} compute the ratio of the decline in log firm premia to the total decline in log hourly wages and find a contribution of 17\%.
 %\footnote{These comparisons are built from \cite{Lachowska2020}, Table 2.} 
 We replicate this exercise in column (7) of Table \ref{tab:evstudytable_body}: loss in premia explains 43\% -> \textcolor{red}{37.1\%} \textcolor{orange}{[OK: Table 2, col. 7 = 0.371]} of hourly wage loss at $t^D+5$. Although our estimates of total hourly wage losses at $t^D+5$ are roughly comparable across the two studies (around -10\%), we find that losses in firm wage premia contribute substantially more to hourly wage losses in our context.
 As we also happen to find larger earnings losses than \cite{Lachowska2020}, the contribution of premium loss to earnings loss at $t^D+5$ that we obtain for the LE sample (14\% -> \textcolor{red}{11.8\%, yay!} \textcolor{orange}{[11.6\% with the rounded Table 2 values (0.037/0.318); 11.8\% only if computed from unrounded coefficients---confirm]}) is close to theirs (12\%).\footnote{We find stronger total earnings and hours effects than \cite{Lachowska2020}. At $t^D+5$, log earnings in the LE sample (respectively, MLO30) decrease by $-0.304$ ($-0.159$) -> \textcolor{red}{-0.316 (-0.159)} \textcolor{orange}{[$-0.318$ ($-0.159$): Table 2, col. 3 and Appendix Table A12, col. 3 at $t^D+5$]} and log hours by $-0.223$ ($-0.090$) -> \textcolor{red}{-0.231 (-0.091)} \textcolor{orange}{[OK: Table 2, col. 4; Appendix Table A12, col. 4]}, as compared to $-0.164$ and $-0.047$ in \cite{Lachowska2020}. This can be explained by three important differences between our setting and theirs: \textit{i)} they identify treated workers using mass separations during the $2008$--$2010$ Great Recession; \textit{ii)} they identify displacements through a $30\%$+ employment drop at the firm level and \textit{iii)} their controls are required to remain at the same firm with continuous employment over the post-event window.}
\cite{Schmieder2020} find a substantially larger contribution of firm premia to daily wage losses, around 76\%.\footnote{Table 3 of \cite{Schmieder2020} provides estimates of losses in log daily wages and firm wage premia: 76.3\% (6.94/9.09) in recessions and 75.5\% (3.91/5.16) in booms. Interestingly, although wage losses differ between booms and recessions, the contribution of firm premia (our main object of interest) appears stable over the cycle.} In a contemporaneous work, \citet{bertheau_etal_2023} find that, across seven European countries, firm premia account for 35 to 100\% of the wage loss. Our 37\% figure for the LE sample lies at the lower end of this range. Our MLO30 estimate (63\%), which relies on a mass-layoff design close to theirs, is near their estimate for France (68\%). Since both studies remove displaced workers' trajectories from the AKM estimation sample, the lower LE figure is unlikely to reflect contamination of firm premia; it is consistent instead with LE workers coming less often from high-premium firms (Appendix Figure~\ref{fig:sample_FEdist}). The remaining gap may also reflect differences in wage concepts (daily versus hourly), in job-loss definition (economic versus mass layoffs), or in data source (exhaustive DADS \textit{postes} versus the 1/12 sample for estimating firm wage premia). They discuss the comparison with our paper in their Appendix C3.

Overall, this analysis shows that losses in firm premia are an important driver of decline in hourly wage in the wake of job displacement, and that reallocation to low-wage firms appears quantitatively relevant in our setting in line with recent work in European context \citep[e.g.][]{Gulyas2019,Schmieder2020}.


\subsection{Robustness tests and additional results}
\label{sec:robustness}

We begin by addressing concerns related to the identification of economic layoffs in small firms. In such firms, firm-level difficulties may partly reflect the performance of individual workers, so that layoffs classified as economic may in practice be related to worker-specific factors. Figure \ref{fig:wage_decomposition_size10} shows that our main results are robust to restricting the sample to firms with (more than) -> \textcolor{red}{at least} \textcolor{orange}{[OK: the caption of the corresponding appendix figure says ``at least 10 employees'']} 10 workers.

We then assess the sensitivity of our results to the construction of AKM firm wage premia. Appendix Table~\ref{tab:paper_R_evstudy_god2} and Figure~\ref{fig:paper_R_wage_decomposition_god2} reproduce our results using an unrestricted AKM specification that does not exclude post-displacement observations from the AKM estimation sample. Consistent with the contamination concerns emphasized in the literature, the unrestricted AKM yields substantially larger estimated losses in firm wage premia, although the overall qualitative patterns remain unchanged.

\paragraph{Time variation in firm wage premia.} We estimate firm wage premia over 2002--2015 and treat them as fixed firm characteristics. Firm pay policies do move over time \citep{lachowska2023drift,engbom2023firmpay}, although \citet{lachowska2023drift} find them to be highly persistent. To check that our results do not hinge on the estimation window, Appendix Figure~\ref{fig:akm_windows} re-estimates the premia over five sub-periods (2002--2005, 2002--2010, 2005--2010, 2005--2015 and 2010--2015) and re-runs the event study with each. The loss in firm wage premia appears in every window, with the same timing. At $t^D+5$, it ranges from $-0.027$, with premia estimated over 2002--2005, a window that largely predates the displacement events, to $-0.046$, with premia estimated over 2005--2010, against $-0.037$ in the baseline.

Next, we replicate our analysis using the mass-layoff sample (MLO30) in Figure \ref{fig:fe_current_event_mlo_pretrends}. We find qualitatively similar patterns in terms of earnings and hourly wage losses. At \(t^D+5\), hourly wage losses are \(-0.070\) in the MLO30 sample versus \textcolor{orange}{\(-0.101\)} in the LE sample. Losses in firm wage premia explain around \textcolor{orange}{\(63\%\)} of hourly wage losses at \(t^D+5\) in MLO30 (versus \textcolor{orange}{\(37\%\)} in LE), a reversal that reflects the over-representation of larger and higher-premium firms in the mass-layoff design. %The MLO sample also provides a longer observation window, allowing us to verify that losses in firm wage premia persist up to \(t^D+8\).

These differences across samples are consistent with differences in sample composition. Displaced workers in MLO30 tend to originate from larger and higher-premium firms than those in LE, in line with the over-representation of large firms induced by the mass-layoff design. Appendix Figure~\ref{fig:sample_FEdist} reports the distribution of pre-displacement firm fixed effects in our two samples and across previous studies, and suggests that the LE sample is somewhat less skewed toward high-premium firms.\footnote{In the top two quintiles of the AKM distribution, the share of displaced workers is \(55.9\%\) in LE and \(67.7\%\) in MLO30 \textcolor{orange}{[53.9\% and 65.8\%, as printed on the R2 version of Appendix Figure~\ref{fig:sample_FEdist}; the 67.7\% of \citet{bertheau_etal_2023} below is unchanged]}; the corresponding shares are \(67.7\%\) in \citet{bertheau_etal_2023}, \(61.8\%\) in \citet{Lachowska2020}, and \(75.9\)--\(76.8\%\) in \citet{Schmieder2020} (recession and boom).}

Finally, we re-estimate the wage decomposition separately for men and women in the LE sample (Appendix Figure~\ref{fig:fe_current_event_by_sex}). The two profiles are quantitatively similar across all four outcomes: at \(t^D+5\), the hourly-wage scar is \(-0.105\) for men versus \(-0.089\) for women, while the AKM firm-premium scar is \(-0.041\) versus \(-0.045\). The reallocation patterns we document are therefore not driven by either sex. (\textcolor{red}{note to me (pw): numbers to swap, including at figure A7}) \textcolor{orange}{[cannot be checked: the by-sex figure prints no coefficients and no export carries them]}





\iffalse
COMPARISON WITH DECOMPOSITION IN LEVEL
\begin{itemize}
	\item Larger effect in log than level
	\item Also wage effect do not dominate in log at the end of the period unlike with the decomposition in level
	\item However the pattern of increasing importance effect of wage effect over hours effect is the same across analysis. 
\end{itemize}
\fi



\FloatBarrier

